\subsection{Pooled Neural Network Prediction}

% ---- Experiment 7: PWDrivers (NN Forecasting)

In Experiment 7, we evaluate pooled country-year prediction of plastic waste generation using a multilayer perceptron (MLP) on the \texttt{PWDrivers} dataset. The data are randomly partitioned into training and held-out test sets after shuffling; therefore, this experiment evaluates prediction across heterogeneous country-year records rather than temporal forecasting from earlier to later observations. We remove outliers from the training set using an isolation forest with a contamination rate of \(r_c = 0.1\). The isolation forest is fitted only on training-partition predictor variables, and test records do not participate in threshold determination. Because records from the same country may appear in both partitions, this experiment does not evaluate generalization to unseen countries or future years. Policy-related factors may influence plastic waste generation but are not represented in the available predictor set, so this experiment does not test whether the MLP captures policy-related effects.

\begin{table}[!ht]
	\centering
	\caption{\centering Neural-network forecasting performance for plastic waste generation on the \texttt{PWDrivers} dataset.}
	\begin{tabular}{@{}lccccc@{}}
		\toprule
		\textbf{Experiment} & \textbf{MAE}             & \textbf{RMSE}         & $\mathbf{R^2}$ & \textbf{WAPE} & \textbf{MAPE} \\
		\midrule
		NN Forecasting      & \(9.146 \times 10^4\)    & \(1.495 \times 10^5\) & 0.726          & 41.22\%       & 92.91\%       \\
		\hline
	\end{tabular}
	\label{tab:nn_forecasting_metrics}
\end{table}


Table~\ref{tab:nn_forecasting_metrics} reports the mean WAPE and \(R^2\) across 16 runs. The results indicate that the MLP captures only broad predictive structure and that its predictions remain far from the observed values. This is consistent with the heterogeneity of the pooled \texttt{PWDrivers} records, which are drawn from different organizations and methodologies. We therefore treat Experiment 7 as an exploratory case study rather than a deployable prediction model; it provides no evidence about unobserved policy-related effects.

% ---- Experiment 8: PWDrivers (Japan & Slovenia); Multivariate Forecasting ----
% Omitted from this manuscript. The source is retained for potential future use.
%
% In Experiment 8, we evaluate conditional prediction of plastic waste generation using GPR, SVR, and MLP models on the \texttt{PWDrivers (Japan)} and \texttt{PWDrivers (Slovenia)} datasets. For each held-out country-year, the models receive the GDP and population variables from that same year, as well as the year itself, to predict plastic waste generation. Thus, this experiment is not a fully autonomous future-forecasting task: deployment would require those contemporaneous covariates to be observed or forecast separately. For the MLP, we first pretrain the model on pooled \texttt{PWDrivers} data after excluding all rows from the target country and removing outliers with an isolation forest using a contamination rate of \(r_c = 0.1\). We then train it on the target country's training data. This experiment compares conditional-prediction performance across methods on tiny datasets. For the \texttt{PWDrivers (Slovenia)} dataset, which contains only \(6\) samples, hyperparameter tuning is not performed due to the limited data availability.
% In Experiment 8, we evaluate conditional prediction of plastic waste generation using GPR, SVR, and MLP models on the \texttt{PWDrivers (Japan)} and \texttt{PWDrivers (Slovenia)} datasets. For each target country, we first partition its country-year records into training and held-out test observations. For each held-out country-year, the models receive the GDP and population variables from that same year, as well as the year itself, to predict plastic waste generation. Thus, this experiment is not a fully autonomous future-forecasting task: deployment would require those contemporaneous covariates to be observed or forecast separately. For the MLP, we first pretrain the model on pooled \texttt{PWDrivers} data after excluding all records from the target country and removing outliers with an isolation forest using a contamination rate of \(r_c = 0.1\). Consequently, neither target-country training observations nor held-out test observations are used during pretraining. We then fine-tune the pretrained model using only the target country's training partition and evaluate it on the held-out observations. This experiment compares conditional-prediction performance across methods on tiny datasets. For the \texttt{PWDrivers (Slovenia)} dataset, which contains only \(6\) samples, hyperparameter tuning is not performed due to the limited data availability.
%
% \begin{table}[!ht]
	\centering
	\caption{\centering Multivariate forecasting performance for plastic waste generation in Japan using the \texttt{PWDrivers (Japan)} dataset.}
	\begin{tabular}{@{}lccccc@{}}
		\toprule
		\textbf{Experiment}         & \textbf{MAE}                   & \textbf{RMSE}                  & \textbf{WAPE}       \\
		\midrule
		GPR                         & \(4.549 \times 10^5\)          & \(4.562 \times 10^5\)          & 11.14\%             \\
		SVR                         & \(4.782 \times 10^5\)          & \(4.792 \times 10^5\)          & 11.71\%             \\
		\textbf{NN Forecasting}     & \(\mathbf{3.943 \times 10^5}\) & \(\mathbf{3.957 \times 10^5}\) & \(\mathbf{9.65}\)\% \\
		\hline
	\end{tabular}
	\label{tab:multivariate_forecasting_japan_metrics}
\end{table}

% \begin{table}[!ht]
	\centering
	\caption{\centering Multivariate forecasting performance for plastic waste generation in Slovenia using the \texttt{PWDrivers (Slovenia)} dataset.}
	\begin{tabular}{@{}lccccc@{}}
		\toprule
		\textbf{Experiment}                & \textbf{MAE}                   & \textbf{RMSE}                  & \textbf{WAPE}       \\
		\midrule
		NN Forecasting                     & \(7.262 \times 10^3\)          & \(8.274 \times 10^3\)          & 11.28\%             \\
		GPR                                & \(1.090 \times 10^4\)          & \(1.121 \times 10^4\)          & 16.94\%             \\
		\textbf{SVR}                       & \(\mathbf{5.580 \times 10^3}\) & \(\mathbf{7.157 \times 10^3}\) & \(\mathbf{8.67}\)\% \\
		\hline
	\end{tabular}
	\label{tab:multivariate_forecasting_slovenia_metrics}
\end{table}

%
% Table~\ref{tab:multivariate_forecasting_japan_metrics} shows that, on the \texttt{PWDrivers (Japan)} dataset, the pretrained MLP achieves the best performance across all reported metrics, with a WAPE of 5.82\%—roughly half the error of GPR (11.14\%) and SVR (11.67\%), which are close to each other. The MLP's advantage suggests that pretraining on the pooled driver data lets the network learn latent driver--target relationships that transfer usefully to the Japanese series, whereas the two kernel-based models, fitted only on the small target training set, cannot recover them. Because only two observations are held out for testing, however, these results should be interpreted as indicative rather than definitive.
%
% Table~\ref{tab:multivariate_forecasting_slovenia_metrics} shows a different pattern on the \texttt{PWDrivers (Slovenia)} dataset. Here SVR attains the lowest error (WAPE 8.67\%), while the MLP performs worst (17.17\%) and GPR is intermediate (16.94\%). With only six samples and no hyperparameter tuning, the flexible MLP cannot convert its pretraining into an advantage and instead fails to adapt to the target series, whereas the simpler SVR generalizes more robustly. Given the tiny dataset, these rankings are likewise indicative rather than definitive.
%
% Taken together, the conditional-prediction results in Experiment 8 are dataset-dependent: the pretrained MLP has the lowest WAPE for Japan, whereas the simpler SVR has the lowest WAPE for Slovenia. Pretraining on the pooled data appears to help only when the target series is large and regular enough to benefit from the transferred representation; for the smallest, noisiest targets, simple kernel methods remain more reliable. These results do not establish autonomous future-forecasting performance, because the observed same-year covariates are supplied to the models. We omit \(R^2\) from both tables because, on hold-out sets of two (Japan) and a handful (Slovenia) of observations, its denominator is near zero and the values become large and negative, reflecting the test-set size rather than model quality, and WAPE is thus the operative metric for this experiment.
